\begin{lemma}[Splitting at a minimum above the midpoint]\label{fprof:split}
Let \(0<a\le b\le1\), and let \(g:[0,N]\to\R\) be 1-Lipschitz with
\(ax\le g(x)\le bx\). For \(3N/4\le n_0<N\), there is a decreasing admissible
chain from \(n_0\) through \(N/2\) to zero, of length
\(O(1+\log(N/(N-n_0)))\), with cost at most \(N\max\{E(a,b),L(a,b)\}\), where
$$
 E(a,b)=\frac{1-2a-2a^2+(2+a)b}{6(1+a)},
$$
$$
 L(a,b)=\begin{cases}
 \dfrac{1+6b-10a}{8},&b\le1/2,\\[5pt]
 \dfrac{b-a}{1+2b}+\dfrac b2-\dfrac{3a}{4},&b\ge1/2.
 \end{cases}
$$
For a grid profile, add \(O(T)\) on mesh \(T\), and for barriers with additive
error \(eN\), add \(O(KeN+KT)\), where \(K\) is the length bound.
\end{lemma}
\begin{proof}
Scale to \(N=1\). Choose a global minimum \(t\) of \(g\) on \([1/2,1]\), and put
\(m=g(t)\), \(v=g(1/2)\), \(e=g(1)\).

We first record the potential argument in the form needed here. If
\(H(z)=hz+\alpha\) majorizes \(g(z)\) on \([q,1]\), with \(h\ge1/2\), put
$$
 \Lambda(x)=H(2x-1)-g(x),\qquad
 V(x)=\frac{\Lambda(x)_+}{1+2h}.
$$
On the region where \(2x-1\ge q\), \(\Lambda\) is nondecreasing, and on
\(\{\Lambda>0\}\) one has
$$
 (-g')_+\le\frac{2h-g'}{1+2h}=V'.
$$
For a nonnegative derivative the right side is nonnegative; for a negative derivative
the inequality reduces to \(g'\ge-1\). Consequently maximal admissible jumps
through this positive region have total cost at most \(V(1)\).
If \(\Lambda(n)\le0\) and \(m_0=2n-1\ge q\), then
\(g(m_0)\le H(m_0)\le g(n)\). A leftmost minimum on \([m_0,n]\) supplies a
zero-cost admissible jump. This is the construction in the preceding two lemmas.
It may be stopped at a specified threshold even if the positive region has not ended.
Positive-region jumps double complementary depths except at a final clipped jump.
Greedy merging of zero-cost jumps bounds their number by twice a logarithmic count.

\emph{An early minimum.} Suppose \(t\le3/4\). Lipschitz continuity gives
\(g(z)\le H(z)=m+z-t\) for \(z\ge t\). Use the potential with \(h=1\),
\(q=t\); at \(x=(1+t)/2\),
$$
 \Lambda((1+t)/2)=m-g((1+t)/2)\le0.
$$
Thus the construction reaches \(n\le(1+t)/2\), with \(n\ge t\), at cost at most
\((1-t+m-e)/3\). If \(n_0\) was already below this threshold, take no such jumps.
The edge to \(t\) is admissible and has zero cost. The edge from \(t\) to \(1/2\)
is admissible because \(t\le3/4\), and costs \(v-m\). The cost is therefore at most
$$
 v-m+\frac{1-t+m-a}{3}.
$$
Maximize this under the weaker scalar constraints
$$
 m\ge at,\qquad v\le\min\{b/2,m+t-1/2\},\qquad 1/2\le t\le1.
$$
Write \(t_*=(1+b)/(2(1+a))\), which lies in \([1/2,1]\).
For \(t\le t_*\), maximizing over \(m\) gives \(m=b/2-t+1/2\), and the value
\((t+b/2-a)/3\) increases with \(t\).
For \(t\ge t_*\), the maximum occurs at \(m=at\), and the value
\((1-a+3b/2-(1+2a)t)/3\) decreases with \(t\).
Their common value at \(t_*\) is \(E(a,b)\).
Relaxing the location interval from \([1/2,3/4]\) to \([1/2,1]\) only enlarges
this upper bound; no assumption \(t_*\le3/4\) is needed.

\emph{A late minimum.} Suppose \(t>3/4\). Every value on \([1/2,1]\) is at least
\(m\ge at>3a/4\). If \(b\ge1/2\), use the envelope \(H(z)=bz\) and stop
the potential construction at the first \(n\le3/4\). Its cost is at most
\((b-a)/(1+2b)\). If \(b\le1/2\), use instead
\(H(z)=z/2+(b-1/2)/2\) on \([1/2,1]\), as in Lemma~\ref{fprof:affine}.
The cost to that stopping point is at most \((1+2b-4a)/8\).
In both constructions the stopping point is at least \(1/2\): a jump starting
above \(3/4\) ends no lower than \(2n-1>1/2\).
The final edge to \(1/2\) is admissible, and costs at most
$$
 v-m\le\frac b2-\frac{3a}{4}.
$$
If this last upper bound is negative, the late case is impossible since \(m\le v\).
Adding the bounds proves the claimed \(L(a,b)\). This argument also applies
when the global minimum lies beyond \(n_0\): it still bounds every interval minimum.
The last edge \(1/2\to0\) has zero cost in both cases.

For precision concerning endpoints, first perform these constructions for a profile
linear on a grid. Interval minima then occur at grid points. A sign-change crossing
costs at most one mesh, as in Lemma~\ref{fprof:affine}. For the early minimum,
the threshold \((1+t)/2\) may lie halfway between grid points; reaching the lower
grid point before jumping to \(t\) costs at most one further mesh.
All these short edges are admissible for sufficiently fine meshes, since
\(t\le3/4\) and \(n_0<1\). Stop at \(3/4\) on a grid containing that point
in the late case. This gives the stated \(O(T)\) error.
Approximate arbitrary continuous profiles by grid interpolants and use the
bounded-length compactness argument of Lemma~\ref{fprof:affine}. If a minimum's
location switches cases along the approximations, the maximum of the two bounds
still applies. Finally clamp approximate barriers and interpolate; the change
in each edge cost is at most twice the uniform perturbation. Rescaling gives
the result on \([0,N]\).
\end{proof}
